\documentclass[11pt,a4paper]{article}

\usepackage[spanish,es-nodecimaldot]{babel}
\usepackage[T1]{fontenc}
\usepackage[utf8]{inputenc}
\usepackage{lmodern}
\usepackage[margin=2.4cm]{geometry}
\usepackage{amsmath,amssymb}
\usepackage{booktabs}
\usepackage{graphicx}
\usepackage{tikz}
\usepackage{enumitem}
\usepackage{xcolor}
\usepackage{hyperref}
\usepackage{setspace}

\hypersetup{
  colorlinks=true,
  linkcolor=black,
  urlcolor=blue,
  citecolor=black
}
\setstretch{1.08}
\setlist[enumerate]{itemsep=0.25em,topsep=0.4em}
\setlist[itemize]{itemsep=0.25em,topsep=0.4em}
\graphicspath{{../matlab/figuras/}{./}}

\newcommand{\DHL}{\ensuremath{\mathrm{DHL}}}
% Texto que el autor debe completar antes de entregar.
\newcommand{\pendiente}[1]{\textcolor{red}{[#1]}}
% Figura de la simulación: si el PNG aún no existe, deja un recuadro.
\newcommand{\figsim}[2]{%
  \IfFileExists{../matlab/figuras/#1.png}%
    {\includegraphics[width=0.78\linewidth]{#1}}%
    {\fbox{\parbox[c][4cm][c]{0.9\linewidth}{\centering\small
      Ejecutar \texttt{simular\_parcial} en MATLAB: genera\\
      \texttt{matlab/figuras/\detokenize{#1}.png}}}}%
  \caption{#2}}

\begin{document}

\begin{titlepage}
\hypersetup{pageanchor=false}
\centering
\vspace*{0.8cm}
\IfFileExists{logo-uni.png}%
  {\includegraphics[height=3.2cm]{logo-uni.png}\\[0.6cm]}%
  {\pendiente{Logotipo oficial y actual de la UNI: guardar como informe/logo-uni.png}\\[0.6cm]}
{\large Universidad Nacional de Ingeniería}\\[0.4cm]
{\large \pendiente{Nombre oficial de la maestría}}\\[1.6cm]
{\LARGE Trabajo Parcial}\\[0.8cm]
{\Large Robot antropomórfico de 3 grados de libertad:\\[0.2cm]
cinemática, singularidades y control cinemático}\\[2.0cm]
\begin{tabular}{@{}ll@{}}
Autor(es) & \pendiente{Apellidos y nombres} \\[0.35cm]
Curso & Robótica y Sistemas Autónomos \\[0.35cm]
Docente & Ing.\ Wilson Dennis Reyes Vásquez \\[0.35cm]
Semestre & 2026-2 \\
\end{tabular}
\vfill
\hypersetup{pageanchor=true}
\end{titlepage}
\setcounter{page}{2}

\tableofcontents
\newpage

\section{Introducción}

Los manipuladores antropomórficos son la configuración más común en la robótica industrial. Su cadena de tres articulaciones rotativas (cintura, hombro y codo) permite posicionar la punta en un volumen amplio con una estructura compacta, y sobre ella se montan las muñecas que orientan la herramienta (Siciliano et al., 2009). Antes de pensar en la dinámica o en técnicas de inteligencia artificial, cualquier aplicación de estos brazos necesita resolver bien su cinemática: saber dónde está la punta, qué ángulos la llevan a un punto, cómo se relacionan las velocidades y en qué posturas el robot pierde movilidad.

Este trabajo toma como caso de estudio un robot real: el manipulador de tres grados de libertad con eslabones de rigidez variable de Ramzy et al. (2024). Los autores construyeron el banco de pruebas, publicaron sus parámetros de Denavit--Hartenberg y sus dimensiones, y midieron con cámara la posición de la punta en 26 posturas. Esa medición permite contrastar la simulación con datos experimentales y no solo con las ecuaciones.

El trabajo busca responder tres preguntas:
\begin{enumerate}
\item ¿La cadena de Denavit--Hartenberg y las dimensiones publicadas reproducen las posiciones que el propio artículo reporta?
\item ¿En qué posturas el robot es singular, y qué consecuencias tiene eso para el control?
\item ¿Qué leyes de control cinemático permiten llevar la punta a un punto y seguir una trayectoria con este robot?
\end{enumerate}

Para responderlas se simulan en MATLAB la cinemática directa, la cinemática inversa, la cinemática diferencial, las singularidades, el espacio de trabajo y el control cinemático, en el espacio articular y en el cartesiano. La dinámica del artículo, con resortes y amortiguadores, queda fuera del alcance de esta entrega.

\section{Desarrollo}

\subsection{Modelo del robot}

El artículo modela un antropomórfico de tres rotativas. Para la cinemática de posición, los autores tratan los eslabones flexibles como rígidos y usan la convención de Denavit--Hartenberg (Denavit y Hartenberg, 1955). El Cuadro~\ref{tab:dh} reúne sus parámetros con las longitudes de la Tabla~2 del artículo:
\begin{equation}
l_1=0{,}04~\mathrm{m},\qquad l_2=0{,}29~\mathrm{m},\qquad l_3=0{,}240~\mathrm{m}.
\end{equation}
El eje \(z_0\) es vertical y apunta hacia arriba. No hay desplazamiento lateral entre los eslabones (\(d_2=d_3=0\)).

\begin{table}[ht]
\centering
\caption{Parámetros de Denavit--Hartenberg (Ramzy et al., 2024, Tabla~1).}
\label{tab:dh}
\begin{tabular}{@{}ccccc@{}}
\toprule
Articulación & \(\theta_i\) & \(d_i\) & \(a_i\) & \(\alpha_i\) \\
\midrule
1 & \(q_1\) & \(l_1=0{,}04\)~m & 0 & \(90^\circ\) \\
2 & \(q_2\) & 0 & \(l_2=0{,}29\)~m & 0 \\
3 & \(q_3\) & 0 & \(l_3=0{,}240\)~m & 0 \\
\bottomrule
\end{tabular}
\end{table}

La misma tabla del artículo trae masas, centros de masa y momentos de inercia, obtenidos en SOLIDWORKS. Esos datos sirven para el modelo dinámico y no intervienen en esta simulación cinemática.

El artículo no publica los límites articulares. Para el espacio de trabajo y para verificar las soluciones de la inversa se adoptaron
\begin{equation}
q_1\in[-180^\circ,180^\circ],\qquad q_2\in[-30^\circ,90^\circ],\qquad q_3\in[-120^\circ,120^\circ].
\end{equation}
Se eligieron porque con ellos el corte \(x=0\) del espacio de trabajo reproduce la forma de herradura de la Figura~5 del artículo. Son un supuesto de este trabajo, no un dato del banco.

\subsection{Cinemática directa}

Cada fila del Cuadro~\ref{tab:dh} da una matriz homogénea
\begin{equation}
\DHL(\theta,d,a,\alpha)=R_z(\theta)\,T(a,0,d)\,R_x(\alpha),
\end{equation}
y la cadena completa es \({}^{0}T_{3}={}^{0}T_{1}\,{}^{1}T_{2}\,{}^{2}T_{3}\) (Craig, 2006; Barrientos et al., 2007). La posición de la punta es la cuarta columna:
\begin{equation}
\begin{aligned}
p_x &= \cos q_1\,(l_2\cos q_2+l_3\cos(q_2+q_3)),\\
p_y &= \sin q_1\,(l_2\cos q_2+l_3\cos(q_2+q_3)),\\
p_z &= l_1+l_2\sin q_2+l_3\sin(q_2+q_3).
\end{aligned}
\label{eq:directa}
\end{equation}
Estas expresiones coinciden con las ecuaciones (6)--(8) del artículo. Para \(q=(0{,}4;\,0{,}7;\,-0{,}8)\)~rad, el producto de matrices y las ecuaciones~\eqref{eq:directa} difieren en cero, dentro de la precisión doble.

Al reconstruir la cadena se detectaron erratas tipográficas en las matrices impresas del artículo. En la matriz (2), la tercera fila de la rotación aparece en ceros y debe ser \([0,\,1,\,0]\), porque \(\sin(90^\circ)=1\). En la matriz (3), el elemento \((3,3)\) aparece como 0 y debe ser 1. En la matriz (5), la última fila aparece como \([0,0,0,0]\) y debe ser \([0,0,0,1]\). Las ecuaciones de posición (6)--(8) sí son correctas, y son las que se usan.

\subsection{Espacio de trabajo}

El espacio de trabajo se obtuvo barriendo las tres articulaciones dentro de sus límites, con 25 valores por articulación, y pasando cada postura por la cinemática directa. El alcance máximo es \(0{,}57\)~m, igual a \(l_1+l_2+l_3\). Con los límites adoptados, \(x\) e \(y\) varían entre \(-0{,}53\) y \(0{,}53\)~m, y \(z\) entre \(-0{,}35\) y \(0{,}57\)~m. La figura~\ref{fig:ws} del anexo muestra la nube tridimensional y el corte \(x=0\), que pueden compararse con las Figuras~4 y~5 del artículo.

\subsection{Cinemática inversa}

Se usó el método geométrico (Barrientos et al., 2007; Reyes Vásquez, 2026). La cintura apunta al punto o al lado opuesto:
\begin{equation}
q_1=\operatorname{atan2}(p_y,p_x)\quad\text{o}\quad q_1=\operatorname{atan2}(p_y,p_x)+\pi .
\end{equation}
Con \(A=p_x\cos q_1+p_y\sin q_1\) (que vale \(\pm\sqrt{p_x^2+p_y^2}\)), el problema se reduce a un brazo de dos eslabones en el plano vertical, con altura \(p_z-l_1\):
\begin{align}
\cos q_3&=\frac{A^2+(p_z-l_1)^2-l_2^2-l_3^2}{2 l_2 l_3}, \qquad
\sin q_3=\pm\sqrt{1-\cos^2 q_3},\\
q_2&=\operatorname{atan2}(p_z-l_1,\,A)-\operatorname{atan2}(l_3\sin q_3,\;l_2+l_3\cos q_3).
\end{align}
Las dos cinturas y los dos codos dan hasta cuatro soluciones. Si \(|\cos q_3|>1\), el punto está fuera de alcance. Si \(p_x=p_y=0\), la punta está sobre el eje de \(q_1\) y esa articulación queda indeterminada.

El Cuadro~\ref{tab:ik} muestra las cuatro soluciones para la punta de \(q=(0{,}4;\,0{,}7;\,-0{,}8)\)~rad. Las cuatro reproducen el punto con un error del orden de \(10^{-16}\)~m. Solo dos respetan los límites articulares, y una de ellas es la postura de partida.

\begin{table}[ht]
\centering
\caption{Soluciones de la cinemática inversa para la punta de \(q=(0{,}4;\,0{,}7;\,-0{,}8)\)~rad.}
\label{tab:ik}
\begin{tabular}{@{}lrrrl@{}}
\toprule
Solución & \(q_1\) (rad) & \(q_2\) (rad) & \(q_3\) (rad) & Límites \\
\midrule
Cintura +, codo + & \(0{,}4000\) & \(-0{,}0203\) & \(0{,}8000\) & dentro \\
Cintura +, codo \(-\) & \(0{,}4000\) & \(0{,}7000\) & \(-0{,}8000\) & dentro \\
Cintura \(-\), codo + & \(-2{,}7416\) & \(2{,}4416\) & \(0{,}8000\) & fuera \\
Cintura \(-\), codo \(-\) & \(-2{,}7416\) & \(-3{,}1213\) & \(-0{,}8000\) & fuera \\
\bottomrule
\end{tabular}
\end{table}

El artículo presenta otra inversa, sus ecuaciones (9)--(12), tomada de trabajos previos. Esa formulación no es consistente con su propia directa. En la ecuación (11), el numerador aparece como \(x^2+y^2+z^2-l_2^2\,l_3^2\), con un producto donde corresponde una resta. En el punto de prueba, el argumento del \(\cos^{-1}\) vale \(1{,}785\) y la función no está definida. Aun corrigiendo esa errata, la ecuación no resta \(l_1\) a la altura \(z\), y el ángulo que asigna a \(\theta_2\) es el del codo, es decir, \(\theta_3\). Por eso se usó la inversa deducida directamente de~\eqref{eq:directa}.

\subsection{Cinemática diferencial}

El jacobiano de posición relaciona las velocidades articulares con la velocidad lineal de la punta, \(\dot p=J(q)\,\dot q\) (Siciliano et al., 2009). Derivando~\eqref{eq:directa}:
\begin{equation}
J=\begin{bmatrix}
-\sin q_1\, A & -\cos q_1\, B & -\cos q_1\, l_3\sin(q_2+q_3)\\
\cos q_1\, A & -\sin q_1\, B & -\sin q_1\, l_3\sin(q_2+q_3)\\
0 & A & l_3\cos(q_2+q_3)
\end{bmatrix},
\end{equation}
con \(A=l_2\cos q_2+l_3\cos(q_2+q_3)\) y \(B=l_2\sin q_2+l_3\sin(q_2+q_3)\).

Se hicieron dos comprobaciones:
\begin{itemize}
\item El jacobiano analítico se comparó con uno numérico por diferencias centrales en \(q=(0{,}4;\,0{,}7;\,-0{,}8)\). La máxima diferencia fue del orden de \(10^{-11}\).
\item Sobre una trayectoria articular quíntica de 3~s, la velocidad \(J(q)\,\dot q\) se comparó con la derivada numérica de la posición de la punta. La máxima diferencia fue del orden de \(10^{-5}\)~m/s, para velocidades de hasta \(0{,}25\)~m/s. Esa diferencia es el error de truncamiento de la derivada numérica (figura~\ref{fig:dif}).
\end{itemize}

\subsection{Singularidades}

Una configuración es singular cuando el determinante del jacobiano se anula: el robot pierde un grado de movilidad y algunas velocidades de la punta exigen velocidades articulares infinitas (Barrientos et al., 2007; Siciliano et al., 2009). Para este robot,
\begin{equation}
\det J=-A\, l_2\, l_3\sin q_3 .
\end{equation}
El resultado coincide con el determinante numérico (\(0{,}0230\) en la postura de prueba). Se anula en dos casos:
\begin{itemize}
\item \(\sin q_3=0\): el codo está completamente estirado (\(q_3=0\), en la frontera exterior del espacio de trabajo) o completamente plegado (\(q_3=\pi\)).
\item \(A=0\): la punta queda sobre el eje de la primera articulación. Con \(l_2=0{,}29\)~m y \(l_3=0{,}240\)~m, el brazo alcanza ese eje sin tener el codo estirado ni plegado. Por ejemplo, con \(q_2=2\pi/3\) y el \(q_3\) que anula \(A\), el determinante es cero aunque \(\sin q_3\neq 0\).
\end{itemize}

El determinante no depende de \(q_1\). Por eso puede representarse como un mapa sobre el plano \((q_2,q_3)\), cuyas curvas de nivel cero son las singularidades (figura~\ref{fig:sing}). Sobre una trayectoria, \(|\det J|\) mide la manipulabilidad del robot (Yoshikawa, 1985). En la simulación, los dos controladores se detienen si \(|\det J|<10^{-4}\).

\subsection{Control cinemático}

\paragraph{Espacio articular (punto a punto).} Entre dos posturas obtenidas con la inversa se interpola cada articulación con un polinomio que parte y llega en reposo, con \(t_f=3\)~s (Reyes Vásquez, 2026; Craig, 2006). La cúbica anula la velocidad en los extremos:
\begin{equation}
a_0=q_i,\quad a_1=0,\quad a_2=\frac{3(q_f-q_i)}{t_f^2},\quad a_3=-\frac{2(q_f-q_i)}{t_f^3}.
\end{equation}
La quíntica anula además la aceleración:
\begin{equation}
a_3=\frac{10(q_f-q_i)}{t_f^3},\quad a_4=-\frac{15(q_f-q_i)}{t_f^4},\quad a_5=\frac{6(q_f-q_i)}{t_f^5}.
\end{equation}
La trayectoria de la punta se obtiene pasando cada \(q(t)\) por la directa (figura~\ref{fig:cub}).

\paragraph{Espacio cartesiano.} Se usan leyes basadas en la pseudoinversa del jacobiano \(J^{+}\), integradas con Euler, \(q_{k+1}=q_k+t_s\,\dot q_k\), con \(t_s=0{,}05\)~s (Kelly y Santibáñez, 2003; Reyes Vásquez, 2026):
\begin{align}
\text{Posición:}\quad &\dot q = J^{+}\,K\tanh(e),\qquad K=8I,\\
\text{Seguimiento:}\quad &\dot q = J^{+}\,(\dot h_d+K e),\qquad K=I,
\end{align}
con \(e=h_d-h(q)\). La saturación con \(\tanh\) limita la velocidad cuando el error es grande.

En el control de posición, la punta parte de \(q=(0{,}15;\,0{,}55;\,-0{,}4)\)~rad y debe llegar al punto \((0{,}406;\,0{,}172;\,0{,}272)\)~m. El error final quedó en el orden de \(10^{-16}\)~m. En el seguimiento, la referencia es una curva tridimensional alrededor de la punta de \(q=(0{,}5;\,0{,}9;\,-0{,}8)\)~rad, con amplitudes de \(0{,}03\), \(0{,}03\) y \(0{,}02\)~m. El error máximo fue de \(1{,}4\times 10^{-4}\)~m y el error final de \(6{,}1\times 10^{-5}\)~m (figuras~\ref{fig:err}--\ref{fig:seg}).

\subsection{Validación con los datos experimentales del artículo}

La Tabla~3 del artículo reporta 26 posturas, la posición de la punta medida con cámara y la posición que los autores calcularon con su modelo. Ese conjunto se usó para validar la cinemática directa de este trabajo. Antes se corrigieron dos erratas de la tabla:
\begin{itemize}
\item En las filas 19 a 26 se imprime \(q_3=6{,}6^\circ\). Con ese valor, la fila 26 da \(z=84{,}9\)~mm; con \(q_3=-6{,}6^\circ\) da \(29{,}8\)~mm, y el artículo reporta \(29{,}92\)~mm. Se usó \(-6{,}6^\circ\). En \(x\) e \(y\) el signo no influye.
\item En la fila 26, la \(x\) teórica aparece como 0. Con \(q_1=88{,}69^\circ\) vale unos 12~mm, por lo que esa celda se excluyó.
\end{itemize}

\begin{table}[ht]
\centering
\caption{Error absoluto medio por eje contra la Tabla~3 del artículo (mm).}
\label{tab:val}
\begin{tabular}{@{}lrrrr@{}}
\toprule
Comparación & \(|e_x|\) & \(|e_y|\) & \(|e_z|\) & máx. \\
\midrule
Modelo (Tabla~2) vs.\ experimental & 1{,}00 & 1{,}73 & 3{,}49 & 9{,}93 \\
Modelo del artículo vs.\ experimental & 2{,}91 & 2{,}29 & 2{,}28 & 6{,}98 \\
Modelo (Tabla~2) vs.\ teórico del artículo & 3{,}73 & 3{,}33 & 3{,}06 & 5{,}97 \\
Modelo ajustado vs.\ teórico del artículo & 0{,}34 & 0{,}22 & 0{,}13 & 0{,}56 \\
\bottomrule
\end{tabular}
\end{table}

El Cuadro~\ref{tab:val} resume los resultados. Con las longitudes de la Tabla~2, el modelo reproduce la medición con errores medios de 1 a 3,5~mm, del mismo orden que los que el artículo reporta para su propio modelo (2,90; 2,29 y 2,40~mm).

En cambio, la columna teórica del artículo no se obtiene con \(l_2=0{,}29\)~m. La posición es lineal en las longitudes, así que se ajustaron \(l_1\), \(l_2\) y \(l_3\) por mínimos cuadrados a esa columna. El resultado fue \(l_1=0{,}0398\)~m, \(l_2=0{,}2954\)~m y \(l_3=0{,}2400\)~m. Con esos valores, el error contra la columna teórica baja a menos de 0,6~mm. Todo indica que los autores calcularon sus valores teóricos con \(l_2\approx 0{,}295\)~m, una longitud distinta de la que publican.

\section{Opinión personal}

\pendiente{Esta sección debe redactarla el autor con sus propias palabras, como exige el enunciado.}

\section{Conclusiones}

\begin{itemize}
\item La cadena de Denavit--Hartenberg del artículo, con las longitudes de su Tabla~2, reproduce las ecuaciones de posición (6)--(8). Contra la medición experimental de 26 posturas, da errores medios de 1 a 3,5~mm, comparables a los que reporta el propio artículo. El modelo cinemático publicado es, por lo tanto, válido.
\item La columna teórica de la Tabla~3 del artículo se reproduce con \(l_2\approx 0{,}295\)~m y no con los 0,29~m publicados. Además, la tabla y varias ecuaciones contienen erratas: el signo de \(q_3\) en las filas 19 a 26, las matrices (2), (3) y (5), y la inversa (9)--(12). Verificar el artículo con un modelo propio permitió detectarlas.
\item La inversa geométrica entrega hasta cuatro soluciones exactas. Los límites articulares descartan algunas, y en la práctica la elección entre las restantes depende de la tarea.
\item El robot es singular con el codo estirado o plegado (\(\sin q_3=0\)) y con la punta sobre el eje de la cintura (\(A=0\)). Las leyes basadas en la pseudoinversa funcionan bien lejos de esas posturas, pero deben detectarlas.
\item Tanto el control punto a punto en el espacio articular como el control cartesiano de posición y de seguimiento llevan la punta al objetivo con errores menores a 0,2~mm en este robot.
\end{itemize}

\section{Recomendaciones}

\begin{enumerate}
\item Medir \(l_2\) directamente en el prototipo, o confirmarla con los autores, antes de usar el modelo para la dinámica. Una diferencia de 5~mm basta para explicar los errores sistemáticos entre la Tabla~2 y la Tabla~3.
\item Incluir en los artículos los límites articulares del banco. Sin ellos, el espacio de trabajo solo puede estimarse a partir de figuras.
\item Reemplazar la pseudoinversa por mínimos cuadrados amortiguados cerca de las singularidades \(A=0\) y \(\sin q_3=0\), para mantener acotadas las velocidades articulares en lugar de detener el control.
\item Planificar las trayectorias lejos de la frontera exterior del espacio de trabajo y del eje de la cintura, usando el mapa de \(\det J\) como criterio.
\item Para el trabajo final, extender el modelo a la dinámica con las masas, centros de masa e inercias de la Tabla~2 del artículo, y validarla con su experimento dinámico.
\end{enumerate}

\section*{Referencias}
\addcontentsline{toc}{section}{Referencias}

\begin{list}{}{\leftmargin=1.5em\itemindent=-1.5em\itemsep=0.5em}
\item Barrientos, A., Peñín, L. F., Balaguer, C. y Aracil, R. (2007). \emph{Fundamentos de robótica} (2.ª ed.). McGraw-Hill.
\item Craig, J. J. (2006). \emph{Robótica} (3.ª ed.). Pearson Educación.
\item Denavit, J. y Hartenberg, R. S. (1955). A kinematic notation for lower-pair mechanisms based on matrices. \emph{Journal of Applied Mechanics, 22}(2), 215--221.
\item Kelly, R. y Santibáñez, V. (2003). \emph{Control de movimiento de robots manipuladores}. Pearson Prentice Hall.
\item Ramzy, M. R., Hammad, S. y Maged, S. A. (2024). Kinematic and dynamic modeling of 3DOF variable stiffness links manipulator with experimental validation. \emph{Applied Sciences, 14}(12), 5285. \url{https://doi.org/10.3390/app14125285}
\item Reyes Vásquez, W. D. (2026). \emph{Robótica y sistemas autónomos: material de clase} [Diapositivas y scripts de MATLAB]. Universidad Nacional de Ingeniería.
\item Siciliano, B., Sciavicco, L., Villani, L. y Oriolo, G. (2009). \emph{Robotics: Modelling, planning and control}. Springer. \url{https://doi.org/10.1007/978-1-84628-642-1}
\item Yoshikawa, T. (1985). Manipulability of robotic mechanisms. \emph{The International Journal of Robotics Research, 4}(2), 3--9. \url{https://doi.org/10.1177/027836498500400201}
\end{list}

\appendix
\section{Anexos: esquemas realizados}

\subsection{Esquema de la cadena}

\begin{figure}[ht]
\centering
\shorthandoff{<>}
\begin{tikzpicture}[line cap=round,line join=round]
  \draw[->] (0,0) -- (0,3.1) node[above] {$z_0$};
  \draw[->] (0,0) -- (1.5,0) node[right] {$x_0$};
  \fill (0,0) circle (1.6pt);
  \node[below left] at (0,0) {base};
  \draw[dashed] (0,0) -- (0,0.7);
  \node[left] at (0,0.35) {$l_1$};
  \fill (0,0.7) circle (1.6pt);
  \node[left] at (0,0.7) {$q_1$};
  \draw[very thick] (0,0.7) -- (2.4,1.7);
  \node[above,sloped] at (1.2,1.28) {$l_2$};
  \node at (0.55,0.42) {$q_2$};
  \fill (2.4,1.7) circle (1.6pt);
  \node[below right] at (2.4,1.7) {$q_3$};
  \draw[very thick] (2.4,1.7) -- (3.7,2.7);
  \node[above,sloped] at (3.05,2.2) {$l_3$};
  \fill (3.7,2.7) circle (1.6pt);
  \node[right] at (3.85,2.7) {punta};
\end{tikzpicture}
\shorthandon{<>}
\caption{Cadena del antropomórfico, vista lateral. \(q_1\) gira alrededor de \(z_0\); \(q_2\) y \(q_3\), alrededor de ejes horizontales paralelos. El dibujo no está a escala.}
\label{fig:cadena}
\end{figure}

\subsection{Figuras de la simulación}

Las figuras siguientes las genera el script \texttt{simular\_parcial.m}, que las guarda en la carpeta \texttt{matlab/figuras/}.

\begin{figure}[htbp]
\centering
\figsim{01_cadena_en_la_postura_de_prueba}{Cinemática directa: base, articulaciones y punta en \(q=(0{,}4;\,0{,}7;\,-0{,}8)\)~rad.}
\label{fig:cadsim}
\end{figure}

\begin{figure}[htbp]
\centering
\figsim{02_espacio_de_trabajo}{Espacio de trabajo 3D y corte \(x=0\), para comparar con las Figuras~4 y~5 de Ramzy et al. (2024).}
\label{fig:ws}
\end{figure}

\begin{figure}[htbp]
\centering
\figsim{03_singularidades}{Mapa de \(\det J(q_2,q_3)\), con las singularidades en rojo, y manipulabilidad sobre la trayectoria quíntica.}
\label{fig:sing}
\end{figure}

\begin{figure}[htbp]
\centering
\figsim{04_cinematica_diferencial}{Velocidad de la punta: \(J(q)\,\dot q\) en color y derivada numérica de la posición en negro discontinuo; las curvas se superponen.}
\label{fig:dif}
\end{figure}

\begin{figure}[htbp]
\centering
\figsim{05_validacion_tabla_3}{Validación contra la Tabla~3 del artículo: modelo frente a medición, y error por postura.}
\label{fig:val}
\end{figure}

\begin{figure}[htbp]
\centering
\figsim{06_error_de_control}{Norma del error en el control de posición y en el seguimiento.}
\label{fig:err}
\end{figure}

\begin{figure}[htbp]
\centering
\figsim{07_control_de_posicion}{Control de posición: cadena inicial, cadena final y camino de la punta.}
\label{fig:pos}
\end{figure}

\begin{figure}[htbp]
\centering
\figsim{08_seguimiento}{Seguimiento cartesiano: referencia y trayectoria de la punta.}
\label{fig:seg}
\end{figure}

\begin{figure}[htbp]
\centering
\figsim{09_interpolador_cubico}{Punto a punto articular: \(q(t)\) del polinomio cúbico y trayectoria resultante de la punta.}
\label{fig:cub}
\end{figure}

\end{document}
